% ch1_a 狭义相对论与平直时空 (书页 1-12 / p016-p027)
% 第1章 Special Relativity and Flat Spacetime
\chapter{狭义相对论与平直时空}

% 1.1 Prelude
\section{引子}

广义相对论（general relativity，GR）是 Einstein 关于空间、时间和引力的理论。究其本质，这是一个非常简单的学科（比方说，与任何涉及量子力学的东西相比）。其核心思想再直截了当不过：自然界中的大多数力都用定义在时空上的场来表示（比如电磁场，或者亚核力所特有的短程场），而引力却内禀于时空自身之中。特别地，我们所体验到的“引力”，正是时空\emph{曲率}（curvature）的一种表现。

这样看来，我们的任务是清楚的。我们需要理解时空，需要理解曲率，还需要理解曲率如何变成引力。大致说来，本书前两章致力于探讨时空，第三章讨论曲率，第四章阐释曲率与引力之间的关系，然后才进入这个理论的各种应用。不过，让我们先放纵自己一下，对将要到来的内容做一番简短的预览，这或许会为我们旅程的头几步提供一点动力。

GR是关于引力的理论，所以不妨从回忆我们先前拥有的引力理论——Newton的理论——入手。它有两个基本要素：一个描述引力场如何受物质影响的方程，和一个描述物质如何响应这个场的方程。这些规则的Newton式传统表述用的是粒子之间的力：质量为$M$和$m$、由矢量$\mathbf{r}=r\,\mathbf{e}_{(r)}$相隔的两个物体之间的力，就是著名的平方反比定律，

\begin{equation*}
\mathbf{F} = -\frac{GMm}{r^2}\,\mathbf{e}_{(r)},
\tag{1.1}
\end{equation*}

这个力作用在质量为$m$的粒子上，按照Newton第二定律使它获得加速度，

\begin{equation*}
\mathbf{F} = m\mathbf{a}.
\tag{1.2}
\end{equation*}

等价地，我们也可以使用引力势$\Phi$的语言；势与质量密度$\rho$由Poisson方程相联系，

\begin{equation*}
\nabla^2\Phi = 4\pi G\rho,
\tag{1.3}
\end{equation*}

而加速度由势的梯度给出，

\begin{equation*}
\mathbf{a} = -\nabla\Phi.
\tag{1.4}
\end{equation*}

式 (1.1) 与式 (1.2)，或者式 (1.3) 与式 (1.4)，都足以定义Newton引力。而要定义GR，我们需要把它们中的每一组都替换成关于时空曲率的表述。

困难的部分是那个支配时空曲率对物质与能量之响应的方程。我们最终会以 Einstein 方程的形式找到想要的东西，

\begin{equation*}
R_{\mu\nu} - \frac{1}{2}R\,g_{\mu\nu} = 8\pi G\,T_{\mu\nu}.
\tag{1.5}
\end{equation*}

这个方程看上去比它本应有的样子更令人生畏，这多半得怪那些希腊下标。事实上，它只不过是$4\times4$矩阵之间的一个方程，下标标记的是每个矩阵的元素。左边的表达式量度时空的曲率，右边的则量度物质的能量和动量，所以正如所许诺的，这个方程把能量与曲率联系了起来。不过，对 Einstein 方程内部运作的详细理解，我们要留到以后再讲。

物质对时空曲率的响应则稍微容易把握一些：自由粒子沿着“可能的最短距离”的路径，也就是测地线（geodesic）运动。换句话说，粒子尽力沿直线运动，但在弯曲时空中可能根本不存在（我们在欧几里得几何中所熟悉的那种意义上的）直线，于是它们只能退而求其次。它们的参数化路径$x^\mu(\lambda)$服从测地线方程（geodesic equation）：

\begin{equation*}
\frac{d^2x^\mu}{d\lambda^2} + \Gamma^\mu_{\rho\sigma}\frac{dx^\rho}{d\lambda}\frac{dx^\sigma}{d\lambda} = 0.
\tag{1.6}
\end{equation*}

到了这一步，你对式 (1.6) 的理解不会比对式 (1.5) 多多少；但用不了多久，这一切都会豁然开朗。

我们后面将会讨论，测地线运动的普适性是GR极其深刻的一个特征。这种普适性正是我们如下论断的根源：引力其实并不是一种“力”，而是时空的一种性质。电场中的带电粒子会感受到一个加速度，使它偏离直线运动；与此相反，引力场中的粒子沿着的是世上最接近直线的东西。这样的粒子感受不到加速度；它们在自由下落。一旦我们更加熟悉GR的精神，下面这个想法就会显得完全合理：在空中飞行的球，比放在桌子上的球更真正地“未被加速”；放在桌上的那个球正在被推离它本来愿意所处的测地线（这就是为什么我们站在地球上时双脚会感受到力）。

我们描述时空曲率所依据的基本概念，将是度规张量（metric tensor），通常记作$g_{\mu\nu}$。度规通过表达对毕达哥拉斯定理$(\Delta l)^2=(\Delta x)^2+(\Delta y)^2$的偏离来编码一个空间的几何（其中$\Delta l$是定义在笛卡尔（Cartesian）网格上的两点之间的距离，两点的坐标间隔为$\Delta x$和$\Delta y$）。这个熟悉的公式只在传统的欧几里得几何中有效，那里隐含地假定了空间是平直的。一旦曲率出现，我们根深蒂固的几何观念就会开始失灵，而我们可以通过追踪毕达哥拉斯关系如何被改变来刻画曲率的大小。这一信息包含在度规张量之中。从度规出发，我们将推导出Riemann曲率张量——既用它来定义 Einstein 方程，也用它得到测地线方程。搭建这套数学装置，就是接下来几章的主题。

尽管为了定量地讨论曲率，不得不引入一定数量的形式语言，GR的核心观念（“引力就是时空的曲率”）却是相当简单的。那么，为什么GR至少在某些蒙昧的圈子里落得了困难、甚至玄奥的名声？因为 Einstein 理论的优雅真理被某些前相对论观念的层层累积所遮蔽；这些观念尽管非常有用，却必须先被抛弃，我们才能按GR的面貌领会这个世界。具体说来，我们生活在一个时空曲率非常小、而且粒子相对于光速大多运动得相当缓慢的世界里。因此，Galileo和Newton的力学对我们来说来得十分自然，尽管它只是那个更深的故事的一种近似。

所以，我们将着手学习这个更深的故事，办法是逐渐剥去那些有用却误导人的Newton直觉的层层外皮。第一步——也就是本章的主题——是探讨狭义相对论（special relativity，SR），即不存在引力（曲率）时的时空理论。希望这部分内容对读者来说大部分是复习，因为我们接下来会走得有些快。要点一方面在于重温SR究竟讲的是什么，另一方面在于引进张量和若干随后至关重要的相关概念，而不必同时承受曲率压在其他一切东西之上的额外复杂性。因此，本章我们将始终在平直时空中工作，而且只使用惯性（类似笛卡尔的）坐标。不用说，在你喜欢的任何坐标系里做SR都是可能的，但事实表明，为此引进必要的工具本身就会把我们带到通往弯曲空间的半路上，所以我们把这件事推迟一段时间。

% 1.2 Space and Time, Separately and Together
\section{空间与时间：分开来看与合起来看}

一种纯粹冷血的GR讲授方式会把第2章（流形）和第1章（狭义相对论与平直时空）的顺序颠倒过来。\emph{流形}（manifold）是用来描述时空的那类数学结构，而\emph{狭义相对论}则是一个调用某种特定时空（没有曲率、因而没有引力的时空）的模型。不过，既然你在读这本书，大概对狭义相对论（SR）至少有些熟悉，对流形却可能一无所知。所以我们的第一步是探访相对熟悉的SR领地，并利用这个机会引入一些对后续发展至关重要的概念和记号。

狭义相对论是关于时空结构的理论，而时空是粒子和场在其中演化的背景。SR充当了Newton力学的替代品，后者同样是一个关于时空结构的理论。无论哪种情形，我们都可以把这个基本结构与支配具体系统的各种动力学定律区分开来：Newton引力是置于Newton力学框架之内的一个动力学系统的例子，而麦克斯韦的电磁学则是在狭义相对论框架内运作的一个动力学系统。

%FIG{1.1}: {在Newton时空中存在一种绝对的切片，把不同时刻的空间切成各不相同的拷贝。粒子世界线被约束只能随时间向前运动，但能以任意速度穿越空间；对于空间中不同地点的两个事件是否发生在同一时刻这个问题，存在人人一致的答案。}

\textbf{时空}（spacetime）是一个四维集合，其元素由三个空间维度和一个时间维度来标记。（下一章我们会更严格地处理这些定义。）时空中的单独一个点称为一个\textbf{事件}（event）。粒子的路径是穿过时空的一条曲线，即一个参数化的一维事件集合，称为\textbf{世界线}（worldline）。这样的描述对SR和Newton力学同样适用。无论哪种情形，“时间”所受的对待显然与“空间”有些不同；具体地说，粒子总是沿时间向前行进，而在空间中却可以自由地来回移动。

然而，SR中粒子可以走的路径集合与Newton理论中的相比，有一个重要差别。在Newton力学中，时空有一个基本划分，分成轮廓分明的“固定时刻的全部空间”切片。\emph{同时性}（simultaneity）的观念——两个事件发生在同一时刻——是毫不含糊地定义了的。粒子的轨迹在时间上永远向前，但除此之外不受约束；特别是，两个这类粒子之间的相对速度没有上限。

在SR中，情况发生了戏剧性的改变：特别是，\emph{两个分离的事件“在同时”发生，这一观念没有良定义。}这并不是说时空完全没有什么结构。毋宁说，在每个事件处我们都可以定义一个\textbf{光锥}（light cone），它是穿过时空的那些路径的轨迹，这些路径可能被穿过该事件的光线所取。Newton力学中把时空唯一地切成由时间参数化的空间切片的那种绝对划分，被一条规则取代：物理粒子不能跑得比光快，因而其运动路径始终留在这些光锥之内。

%FIG{1.2}: {在狭义相对论中，不存在"某一时刻的全部空间"这种绝对观念。取而代之的是一条规则：粒子永远以不大于光速的速率行进。于是我们可以在每个事件处定义光锥，它们局域地刻画了允许轨迹的集合。对于彼此的光锥之外的两个事件，哪个事件在时间上先发生，没有普适的观念。}

SR中不存在偏好的时间切分，这正是时空观念在此语境中比在Newton力学中更为根本的关键所在。当然，我们可以在时空中选取具体的坐标系，而且一旦选定，在这个特定系统中谈论分离的事件取相同的时间坐标值就是有意义的；但还会存在其他可能的坐标系，它们通过把空间和时间相互“旋转”而与第一组相联系。这一现象是欧几里得几何中旋转的自然推广，我们现在就转向它。

考虑一块最寻常不过的二维平面。给这种平面上的点做标记，通常方便的做法是引入坐标，例如定义正交的$x$轴和$y$轴，按通常方式把每个点投影到这两个轴上。然而，显然关于这个平面的大多数有趣的几何事实都与我们坐标的选取无关；不存在任何偏好的方向。作为一个简单的例子，我们可以考虑两点之间的距离，它由下式给出：

\begin{equation*}
(\Delta s)^2 = (\Delta x)^2 + (\Delta y)^2.
\tag{1.7}
\end{equation*}

在一个不同的笛卡尔坐标系——由相对原来的轴旋转过的$x'$轴和$y'$轴定义——之中，距离的公式丝毫未变：

\begin{equation*}
(\Delta s)^2 = (\Delta x')^2 + (\Delta y')^2.
\tag{1.8}
\end{equation*}

因此我们说，距离在这样的坐标变更下不变。

%FIG{1.3}: {二维欧几里得空间，画有两个不同的坐标系。"两点之间的距离"这类观念与所选的坐标系无关。}

这就是为什么把平面看作一个内禀的二维空间是有用的，而不是把它看作两个被任意拼到一起的、从根本上不同的一维空间：虽然我们用两个不同的数来标记每个点，这些数并不是几何的本质，因为我们可以在保持距离不变的同时把轴相互旋转。在Newton物理学中，空间和时间却不是这样；把空间和时间相互旋转没有什么有用的含义。相反，“单一时刻的全部空间”这一观念具有独立于坐标的意义。

SR则是另一回事。让我们考虑时空上的坐标$(t,x,y,z)$，按如下方式建立。空间坐标$(x,y,z)$构成一个标准的笛卡尔系，例如可以把相交成直角的刚性杆焊接在一起来构造。这些杆必须自由运动、不受加速。时间坐标由一组时钟定义，这些时钟相对于空间坐标不运动。（既然这是思想实验，我们可以想象这些杆无限长，而且空间中每个点都有一只时钟。）这些时钟按下述意义同步。想象我们从空间中的点1沿直线以恒定速度$c$发送一束光到点2，然后立即回到1（速度为$-c$）。那么，光束到达点2时坐标钟上的时间（记作$t_2$），应当介于光束离开点1时坐标钟上的时间（$t_1$）与它回到同一只钟时的时间（$t_1'$）的正中间：

\begin{equation*}
t_2 = \frac{1}{2}(t_1' + t_1).
\tag{1.9}
\end{equation*}

如此构造出来的坐标系称为一个\textbf{惯性系}（inertial frame），或简称“惯性坐标”。这些坐标是空间中笛卡尔（正交归一）坐标向时空的自然推广。（如此小心翼翼地构造，是为了我们只在\emph{局域}上作比较；例如，绝不在同一时刻去比较两只相距遥远的钟。等到进入广义相对论，这种谨慎将变得更加必要，因为在那里将没有任何办法在整个时空上构造惯性坐标。）

%FIG{1.4}: {在惯性坐标系中校准时钟。当一束光从点1射向点2再反弹回来时，如果时间$t_2$恰好位于$t_1$与$t_1'$的正中间，这些钟就是同步的。}

我们可以用这套流程构造任意多个惯性系，它们与第一个的差别在于初始位置和时间的偏移、角度以及（恒定的）速度。在Newton式的世界里，新坐标$(t',x',y',z')$会有这样的特点：$t'=t+\text{常数}$，与空间坐标无关。也就是说，“两个事件同时发生，即在同一时刻发生”有一个绝对的含义。但在SR中这并不成立；一般而言，由$t=\text{常数}$定义的三维“空间”将不同于由$t'=\text{常数}$定义的那些。

不过，我们并没有完全坠入混乱。暂且不加任何动机，考虑一下我们将称之为两个事件之间的\textbf{时空间隔}（spacetime interval）的东西：

\begin{equation*}
\boxed{(\Delta s)^2 = -(c\Delta t)^2 + (\Delta x)^2 + (\Delta y)^2 + (\Delta z)^2.}
\tag{1.10}
\end{equation*}

（注意，即使对两个不重合的点，它也可以为正、为负或为零。）这里的$c$是空间与时间之间某个固定的换算因子，也就是说，一个固定的速度。作为经验事实，电磁波恰以这个速度$c$在真空中传播，因此我们把它称作“光速”。然而，重要的并不是光子碰巧以这个速率行进，而是存在这样一个$c$，使得\emph{时空间隔在惯性坐标的变更下不变}。换句话说，如果我们建立一个新的惯性系$(t',x',y',z')$，间隔将具有同样的形式：

\begin{equation*}
(\Delta s)^2 = -(c\Delta t')^2 + (\Delta x')^2 + (\Delta y')^2 + (\Delta z')^2.
\tag{1.11}
\end{equation*}

这正是为什么把SR看作四维时空的理论是有意义的，这个四维时空称为\textbf{Minkowski 空间}（Minkowski space）。（它是四维流形的一个特殊情形，流形我们稍后会详细处理。）我们将会看到，我们已经隐式定义的那些坐标变换，确实在某种意义上把空间和时间相互旋转。不存在“同时事件”的绝对观念；两件事是否发生在同一时刻，取决于所用的坐标。因此，把 Minkowski 空间分割成空间和时间，是我们为自己的目的而做的选择，并不是情形本身固有的东西。

与SR相联系的“佯谬”几乎全部源于对Newton观念的顽固固守：坚持存在唯一的时间坐标，坚持存在“单一时刻的全部空间”。改用时空来思考，而不是把空间与时间合在一起思考，这些佯谬往往就会消失。

让我们引入一些方便的记号。时空上的坐标将用带希腊字母上标指标的字母来表示，指标从0取到3，其中0一般表示时间坐标。于是，

\begin{equation*}
x^{\mu}:\quad
\begin{aligned}
x^0 &= ct \\
x^1 &= x \\
x^2 &= y \\
x^3 &= z.
\end{aligned}
\tag{1.12}
\end{equation*}

（不要开始把这些上标当成幂指数。）此外，为简单起见，我们将选取使下式成立的单位制：

\begin{equation*}
c = 1;
\tag{1.13}
\end{equation*}

因此在后面所有的公式里，我们都将省去因子$c$。经验告诉我们，$c$为$3\times10^8$米每秒；于是，我们是在1秒等于$3\times10^8$米的单位制下工作。有时分别指称$x^\mu$的空间分量和时间分量会很有用，所以我们将用拉丁字母上标单独表示空间分量：

\begin{equation*}
x^{i}:\quad
\begin{aligned}
x^1 &= x \\
x^2 &= y \\
x^3 &= z.
\end{aligned}
\tag{1.14}
\end{equation*}

把时空间隔写成更紧凑的形式也很方便。为此我们引入一个$4\times4$矩阵，即\textbf{度规}（metric），用两个下标书写：

\begin{equation*}
\eta_{\mu\nu} =
\begin{pmatrix}
-1 & 0 & 0 & 0 \\
0 & 1 & 0 & 0 \\
0 & 0 & 1 & 0 \\
0 & 0 & 0 & 1
\end{pmatrix}.
\tag{1.15}
\end{equation*}

（一些文献——尤其是场论教材——用相反的符号定义度规，所以要当心。）于是我们有漂亮的公式

\begin{equation*}
(\Delta s)^2 = \eta_{\mu\nu}\Delta x^\mu \Delta x^\nu.
\tag{1.16}
\end{equation*}

%FIG{1.5}: {画在时空图上的一个光锥。图中标出了与原点类空分离、类光分离和类时分离的点。}

这个公式引入了\textbf{求和约定}（summation convention）：既作为上标又作为下标出现的指标要对之求和。我们把这样的指标称为\textbf{哑指标}（dummy indices）；必须记住，它们对所有可能的值求和，而不取任何特定的值。（情形将永远如此：哑指标严格成对出现，一个“在上”，一个“在下”。这一点以后再谈。）因此，式 (1.16) 的内容与式 (1.10) 完全相同。

\textbf{时空图}（spacetime diagram）是一种极其有用的工具，那就让我们从这个观点来考察 Minkowski 空间。我们可以先把起初的$t$轴与$x$轴画成直角，并隐去$y$轴和$z$轴。（时空图上画出的“直角”并不必然意味着“在时空中正交”，尽管在这个例子里，$t$轴与$x$轴恰好确实正交。）考察以速率$c=1$行进的路径颇有启发，这些路径由$x=\pm t$给出。由以光速行进的直线与单个事件相连的全部点的集合，就是\textbf{光锥}；因为如果想象再纳入一个空间维度，那两条对角线就会补全成一个锥面。光锥天然分为未来光锥与过去光锥；一点$p$的未来与过去光锥内部所有点的集合，称为与$p$\textbf{类时分离}（timelike separated）；光锥之外的点与$p$\textbf{类空分离}（spacelike separated）；而恰好位于锥面上的点，则与$p$\textbf{类光分离}（lightlike or null separated）。回到式 (1.10)，我们看到类时分离的点之间的间隔为负，类空分离的点之间的间隔为正，类光分离的点之间的间隔为零。（间隔定义为$(\Delta s)^2$本身，而不是这个量的平方根。）

间隔对类时曲线（比光慢的粒子实际上将沿之运动）为负，这件事有些恼人，所以我们定义\textbf{固有时}（proper time）$\tau$满足

\begin{equation*}
(\Delta\tau)^2 = -(\Delta s)^2 = -\eta_{\mu\nu}\Delta x^\mu \Delta x^\nu.
\tag{1.17}
\end{equation*}

时空间隔的一个关键性质是：\emph{两个事件之间的固有时，量度的是沿连接这两事件的直线路径运动的观者所经历的时间。}这在非常特殊的情形下最容易看清：两个事件具有相同的空间坐标，只在时间上分离；这对应于在两事件之间旅行的观者相对所用坐标系静止。于是$(\Delta\tau)^2=-\eta_{\mu\nu}\Delta x^\mu\Delta x^\nu=(\Delta t)^2$，故$\Delta\tau=\Delta t$；当然，我们本来就把$t$定义为位于固定空间位置上的钟所量度的时间。但时空间隔在惯性系的变更下不变；因此，两个固定事件之间的固有时 (1.17) 在观者运动着的惯性系中求出的值，与在观者静止的惯性系中求出的值相同。

一个关键的事实是：对更一般的轨迹，固有时与坐标时并不相同（尽管固有时总是沿轨迹运动的观者所携带的钟量度的时间）。考虑连接事件$A$和$C$的两条轨迹：一条是直线，经过标记为$B$的中点；另一条由一位观者走完，他以恒定速度$v=dx/dt$从$A$离开行进到一点$B'$，再以恒定速度$-v$返回，与前者相交于事件$C$。

%FIG{1.6}: {双生子佯谬。沿穿过时空的直线路径$ABC$旅行的行者，将比沿非直线路径$AB'C$旅行的人衰老得更多。既然固有时是穿过时空走过的距离的量度，这应当不足为怪。（唯一可能令人惊讶的是，直线路径竟然是固有时\emph{最长}的那一条；这可以追溯到度规类时分量的负号。）}

选取惯性坐标，使直的轨迹描述一个静止不动的粒子，事件$A$位于坐标$(t,x)=(0,0)$，$C$位于$(\Delta t,0)$。两条路径于是在时空中构成一个等腰三角形；$B$的坐标为$(\frac{1}{2}\Delta t,0)$，$B'$的坐标为$(\frac{1}{2}\Delta t,\Delta x)$，其中$\Delta x=\frac{1}{2}v\Delta t$。显然有$\Delta\tau_{AB}=\frac{1}{2}\Delta t$，但是

\begin{equation*}
\begin{aligned}
\Delta\tau_{AB'} &= \sqrt{\left(\tfrac{1}{2}\Delta t\right)^2 - (\Delta x)^2} \\
&= \tfrac{1}{2}\sqrt{1-v^2}\,\Delta t.
\end{aligned}
\tag{1.18}
\end{equation*}

应当显然有$\Delta\tau_{BC}=\Delta\tau_{AB}$和$\Delta\tau_{B'C}=\Delta\tau_{AB'}$。于是，沿直线从事件$A$旅行到$C$的观者经历的时间流逝为$\Delta\tau_{ABC}=\Delta t$，而走出去又返回的那位经历的却是

\begin{equation*}
\Delta\tau_{AB'C} = \sqrt{1-v^2}\,\Delta t < \Delta t.
\tag{1.19}
\end{equation*}

尽管两位观者在时空中的同一点出发、又在同一点结束，他们衰老的程度却不一样。这就是著名的“双生子佯谬”（twin paradox）——形形色色的误解和蹩脚解释的不幸上演之处。真相很直接：时空中的非直线路径与直线路径有不同的间隔，正如空间中的非直线路径与直线路径有不同的长度。当然，这件事并不像听起来那么平凡；深刻的洞见在于“沿世界线流逝的时间”以何种方式与穿过时空历经的间隔相联系。在Newton式的世界里，坐标$t$代表贯穿全部时空的普适的时间之流；在相对论中，$t$只是一个方便的坐标，时间的流逝取决于你沿哪条路径旅行。一个重要的区别是：非直线路径的固有时\emph{更短}。在空间中，两点之间最短的距离是直线；在时空中，两事件之间最长的固有时属于直的轨迹。

并非所有轨迹都齐整到可以由直线段拼成。在更一般的情形下，引入无穷小间隔，即\textbf{线元}（line element），会很有用：

\begin{equation*}
\boxed{ds^2 = \eta_{\mu\nu}dx^\mu dx^\nu,}
\tag{1.20}
\end{equation*}

其中$dx^\mu$是无穷小的坐标位移。（我们这里的处理相当不正式，不过以后会补正。）从这个定义出发，人们很想取平方根并沿一条路径积分，以得到有限的间隔，但$\int\sqrt{\eta_{\mu\nu}dx^\mu dx^\nu}$究竟应当是什么意思，多少有些不清楚。换一种做法，我们把穿过时空的一条路径看作参数化曲线$x^\mu(\lambda)$。注意，与Newton力学的传统做法不同，参量$\lambda$不一定等同于时间坐标。于是我们可以计算导数$dx^\mu/d\lambda$，并把沿类空曲线（无穷小间隔为类空的曲线）的路径长度写作

\begin{equation*}
\Delta s = \int \sqrt{\eta_{\mu\nu}\frac{dx^\mu}{d\lambda}\frac{dx^\nu}{d\lambda}}\,d\lambda,
\tag{1.21}
\end{equation*}

其中积分沿整条路径进行。对类时路径我们用固有时

\begin{equation*}
\Delta\tau = \int \sqrt{-\eta_{\mu\nu}\frac{dx^\mu}{d\lambda}\frac{dx^\nu}{d\lambda}}\,d\lambda,
\tag{1.22}
\end{equation*}

它将为正。（对类光路径，间隔干脆就是零。）当然，我们可以考虑在一些地方类时、在另一些地方类空的路径，但幸运的是很少有必要，因为物理粒子的路径从不改变其特征（有质量粒子沿类时路径运动，无质量粒子沿类光路径运动）。再说一遍，$\Delta\tau$确实就是沿轨迹运动的观者所量度的时间。

狭义相对论中的\emph{加速度}（acceleration）这一概念名声不佳，但毫无道理。在建立惯性坐标时，我们\emph{当然}小心翼翼地确保了在此类坐标中静止的粒子不受加速。然而，一旦建立了这样的坐标，我们就可以随意考虑物理粒子的任何轨迹，无论加速与否。特别是，“SR处理不了加速轨迹、必须请出广义相对论”的传言毫无事实依据。广义相对论是在有引力之时、当时空变得弯曲之际才变得相关。平直时空中的任何过程都在狭义相对论的框架内得到描述；特别地，像式 (1.22) 这样的表达式是完全普遍的。

% 1.3 Lorentz Transformations
\section{洛伦兹变换}

现在，我们可以在比之前稍抽象一点的层次上考虑时空中的坐标变换。我们感兴趣的是如何把经由上述流程构造的各种惯性系联系起来，并对此给出一个形式化的描述；也就是说，（考虑那些）使间隔 (1.16) 保持不变的坐标系。其中简单的一类是\textbf{平移}（translations），它们只是让坐标平移（在空间或时间上）：

\begin{equation*}
x^\mu \rightarrow x^{\mu'} = \delta^{\mu'}_{\mu}\,(x^\mu + a^\mu),
\tag{1.23}
\end{equation*}

其中$a^\mu$是由四个固定数构成的一组数，而$\delta^{\mu'}_{\mu}$是传统Kronecker delta符号的四维版本：

\begin{equation*}
\delta^{\mu'}_{\mu} =
\begin{cases}
1 & \text{当 } \mu' = \mu, \\
0 & \text{当 } \mu' \neq \mu.
\end{cases}
\tag{1.24}
\end{equation*}

注意，我们把撇号加在指标上，而不是加在$x$上。这样做的理由，等我们开始同矢量和张量打交道之后会变得更加清楚；这种记号提醒我们：几何对象是同一个，只是它的分量是相对于另一个坐标系分解的。平移保持差值$\Delta x^\mu$不变，所以间隔不变并不令人惊讶。其他相关的变换还包括空间的\textbf{旋转}（rotations），以及以一个恒定速度矢量做的偏移，即\textbf{推动}（boosts）；这些都是线性变换，可以用一个（不依赖时空的）矩阵去乘$x^\mu$来描述：

\begin{equation*}
x^{\mu'} = \Lambda^{\mu'}{}_{\nu}x^\nu,
\tag{1.25}
\end{equation*}

或者，用更传统的矩阵记号，

\begin{equation*}
x' = \Lambda x.
\tag{1.26}
\end{equation*}

（我们一般会使用指标，而不是矩阵记号，但此刻我们有意把这里的讨论与通常用矩阵描述的某些其他熟悉概念联系起来。）这些变换并不使差值$\Delta x^\mu$保持不变，而是把它们再乘上矩阵$\Lambda$。什么样的矩阵能使间隔不变？坚持用矩阵记号的话，我们想要的是

\begin{equation*}
\begin{aligned}
(\Delta s)^2 &= (\Delta x)^{\mathrm{T}}\eta(\Delta x) = (\Delta x')^{\mathrm{T}}\eta(\Delta x') \\
&= (\Delta x)^{\mathrm{T}}\Lambda^{\mathrm{T}}\eta\Lambda(\Delta x),
\end{aligned}
\tag{1.27}
\end{equation*}

因此

\begin{equation*}
\eta = \Lambda^{\mathrm{T}}\eta\Lambda,
\tag{1.28}
\end{equation*}

% 本文件至此为止：书页12末句止于式 (1.28) 后的逗号，正文延续到书页13，由下一块（ch1_b）接续。
